% !TEX root = ../main.tex

\chapter{算法设计与实验}

\section{任务定义与研究思路}

本文关注的核心任务是：给定机器人当前视觉观察、短历史观察和目标图像或目标节点图像，生成一段局部 waypoint 轨迹，并将其作为中层控制器的参考输入。沿用 NoMaD 的基本思想\cite{nomad2023}，本文不直接预测单步速度命令，而是学习条件轨迹分布
\begin{equation}
p(\mathbf{A}_t \mid \mathcal{O}_t, I_g),
\end{equation}
其中 $\mathcal{O}_t$ 表示时刻 $t$ 的历史视觉观察，$I_g$ 表示目标图像或从 topomap 中筛选得到的局部目标节点，$\mathbf{A}_t$ 表示未来局部 waypoint 序列。

围绕这一任务，本文在算法层的工作可概括为三部分。第一，基于项目代码准确复现 NoMaD 的训练与推理链路，明确视觉编码器、距离预测头与扩散策略头之间的关系。第二，围绕部署可用性对推理层进行优化，重点研究 DDIM 推理、CFG 引导增强和 TTS 推理时搜索。第三，围绕视觉表征能力设计视觉编码器升级实验，并与基线配置进行对比。需要强调的是，本章关注的是“哪些候选方法值得进入闭环验证”，因此所有结论都首先被限定在离线统计与离线代理指标层面。

\begin{figure}[htbp]
    \centering
    \fbox{\rule{0pt}{55mm}\rule{0.84\linewidth}{0pt}}
    \caption{图位占位：NoMaD 算法研究对象与数据流示意图（待补）}
    \label{fig:algo-overview-placeholder}
\end{figure}

\section{NoMaD 基础模型复现}

\subsection{输入输出形式}

根据当前项目配置，NoMaD 的输入由历史视觉观察与目标图像构成。实验中典型设置为 \texttt{context\_size=3}，即模型同时接收若干历史观察帧与一张目标图像。输出则为固定长度的局部轨迹序列，每个 waypoint 用局部坐标系下的二维位移表示：
\begin{equation}
\mathbf{a}_i = (\Delta x_i, \Delta y_i), \qquad i=1,\ldots,H.
\end{equation}
整段轨迹写为
\begin{equation}
\mathbf{A} = [\mathbf{a}_1,\mathbf{a}_2,\ldots,\mathbf{a}_H].
\end{equation}

这一表示的优点在于，它比直接输出底层速度命令更适合作为平台无关的高层接口，也比全局路径更适合闭环局部控制。对于 Lite3 这样的足式平台，高层输出局部轨迹、中层完成控制映射的方式尤为关键。

\subsection{模型结构与训练目标}

NoMaD 由视觉编码器、距离预测头和扩散策略头三部分组成\cite{nomad2023}。视觉编码器负责将历史观察与目标条件编码为统一特征；距离预测头输出当前观察到目标的相对距离估计；扩散策略头在该条件上生成局部轨迹样本。项目代码显示，距离预测头与扩散策略头共享同一视觉条件空间，因此视觉编码器的变化会同时影响二者。

训练时，总损失由距离回归损失与扩散噪声预测损失组成，可写为
\begin{equation}
\mathcal{L} = \alpha \mathcal{L}_{dist} + (1-\alpha)\mathcal{L}_{diff}.
\end{equation}
其中，$\mathcal{L}_{dist}$ 约束候选目标排序能力，$\mathcal{L}_{diff}$ 约束局部轨迹生成质量。本文对训练代码的检查表明，距离预测头与视觉编码器是直接耦合的，因此更换视觉编码器后，距离预测分支必须同步重新训练，否则距离头与轨迹头将共享不一致的条件分布。

\subsection{推理链路的统一封装}

在推理阶段，系统先编码观察与目标，再在需要时进行候选节点距离评估，最后通过扩散采样生成轨迹。为保证后续实验和系统章节可复用，本文将 \texttt{policy\_config}、\texttt{policy\_checkpoint}、\texttt{scheduler\_kind}、\texttt{ddim\_steps}、\texttt{cfg\_weight}、\texttt{num\_samples} 等参数收口到统一推理模块中。这样做的意义在于：推理优化被限制在高层策略层，不需要同时改动状态机、导航主机和平台接口。

\section{基线实验}

为给后续所有优化实验建立统一参考，本文首先复现了原始 NoMaD 离线推理流程。已有结果表明，原始模型参数量约为 19,049,675，DDPM-10 配置下单次离线推理时间约为 0.1250\,s，对应约 8.0\,Hz 的单次离线推理频率。其主要作用不是证明“原始模型足够好”，而是作为比较基线：后续 DDIM、CFG、TTS 和视觉编码器实验都需要回到这一基线来判断收益与代价。

\begin{table}[htbp]
    \centering
    \caption{NoMaD 基线离线推理结果}
    \label{tab:baseline-inference}
    \begin{tabular}{lll}
        \toprule
        指标 & 含义 & 数值 \\
        \midrule
        模型参数量 & 可学习参数总量 & 19,049,675 \\
        DDPM 推理时间 & 单次完整离线推理耗时 & 0.1250\,s \\
        推理频率 & 每秒可执行离线推理次数 & 8.0\,Hz \\
        轨迹采样规格 & 候选轨迹数 $\times$ waypoint 数 & $8\times8$ \\
        距离预测输出 & 候选目标相对距离估计均值 & 11.6589 \\
        \bottomrule
    \end{tabular}
\end{table}

\section{实验指标与统计原则}

为了避免仅凭单次样本或单一指标做结论，本文采用多指标联合分析。主要指标包括：Latency，用于衡量推理开销；CI95，用于表示均值的 95\% 置信区间；MSE vs DDPM-10，用于衡量与基线采样分布的偏移；Forward Progress，用于刻画轨迹在主前进方向上的推进程度；Lateral Abs，用于衡量横向偏摆幅度；Smoothness，用于衡量轨迹变化的平滑性；Dist Pred Mean，用于观察候选目标距离评估是否随条件空间变化发生系统偏移。

其中，CI95 并不是新的性能指标，而是均值稳定性的统计描述。若某个配置平均值较优但 CI95 很宽，说明其结果在不同 case 间波动较大。Forward Progress 也不能单独解释为“越大越好”，因为过大的前进量可能是激进采样带来的分布偏移，必须结合 MSE、Lateral Abs 和 Smoothness 一起分析。

\section{DDIM 推理加速实验}

\subsection{实验设计}

如第二章所述，DDIM 的作用是在不重新训练模型的情况下减少扩散去噪步数，从而降低推理时延\cite{ddim2020}。本文在统一脚本下比较 \texttt{ddpm\_10}、\texttt{ddim\_10}、\texttt{ddim\_5}、\texttt{ddim\_3}、\texttt{ddim\_2} 和 \texttt{ddim\_1} 六种配置，测试集共包含 24 个 case。该实验的目标是回答两个问题：第一，步数减少后推理能快多少；第二，轨迹分布是否仍保持在可接受范围。

\subsection{实验结果}

\begin{table}[htbp]
    \centering
    \caption{DDIM 推理加速实验结果}
    \label{tab:ddim-results}
    \begin{tabular}{lccccc}
        \toprule
        配置 & Latency/ms & CI95/ms & Speedup & MSE vs DDPM-10 & Forward Progress \\
        \midrule
        ddpm\_10 & 39.9790 & 1.3748 & 1.00$\times$ & 0.0000 & 6.4998 \\
        ddim\_10 & 38.0260 & 0.6034 & 1.05$\times$ & 0.0062 & 6.4686 \\
        ddim\_5 & 19.2581 & 0.4319 & 2.08$\times$ & 0.0048 & 6.5257 \\
        ddim\_3 & 11.6579 & 0.3438 & 3.43$\times$ & 0.0124 & 6.6224 \\
        ddim\_2 & 8.2588 & 0.2779 & 4.84$\times$ & 0.0210 & 6.6471 \\
        ddim\_1 & 4.5092 & 0.1427 & 8.87$\times$ & 3.5168 & 9.2789 \\
        \bottomrule
    \end{tabular}
\end{table}

\subsection{结果分析}

结果表明，DDIM 的确能显著降低推理时延，其中 \texttt{ddim\_2} 已接近 5 倍加速。更关键的是，\texttt{ddim\_2} 在 MSE 与 Forward Progress 上仍保持相对平衡，说明它不是以严重破坏轨迹分布为代价换取速度。相比之下，\texttt{ddim\_1} 虽然最快，但 MSE 大幅升高，说明轨迹已明显偏离基线分布，不适合作为稳定部署方案。

因此，本文将 \texttt{ddim\_2} 视为后续闭环验证中最值得优先考虑的推理配置。它回答了一个重要问题：扩散策略并不必然意味着“质量高但太慢”，通过合理采样调度，NoMaD 可以进入更接近真实部署的时延区间。

\begin{figure}[htbp]
    \centering
    \fbox{\rule{0pt}{52mm}\rule{0.84\linewidth}{0pt}}
    \caption{图位占位：DDIM 加速实验结果图（待补）}
    \label{fig:algo-ddim-placeholder}
\end{figure}

\section{CFG 引导增强实验}

\subsection{实验设计}

CFG 的目标是在推理阶段增强轨迹对目标条件的响应\cite{cfg2022}。本文选取 $w\in\{-1.0,0.0,0.5,1.0,2.0,4.0\}$ 六组配置进行比较。实验关注的问题不是“引导越强越好”，而是“多大引导能够在目标牵引与分布稳定性之间取得平衡”。

\subsection{实验结果与分析}

已有结果表明，适度的 CFG 能提高轨迹对目标方向的响应，但代价是更高的推理时延和潜在的分布偏移。总体上，$w=1.0$ 左右在前进性和稳定性之间表现较为均衡，而过大的 $w$ 会带来更明显的横向偏摆和轨迹漂移。因此，本文将 CFG 定位为可调部署旋钮，而非必须开启的固定默认项。

\begin{table}[htbp]
    \centering
    \caption{CFG 引导实验结果占位}
    \label{tab:cfg-results}
    \begin{tabular}{lcccc}
        \toprule
        配置 & Latency/ms & Shift vs $w=0$ & Forward Progress & Lateral Abs \\
        \midrule
        $w=-1.0$ & 待补 & 待补 & 待补 & 待补 \\
        $w=0.0$ & 待补 & 待补 & 待补 & 待补 \\
        $w=0.5$ & 待补 & 待补 & 待补 & 待补 \\
        $w=1.0$ & 待补 & 待补 & 待补 & 待补 \\
        $w=2.0$ & 待补 & 待补 & 待补 & 待补 \\
        $w=4.0$ & 待补 & 待补 & 待补 & 待补 \\
        \bottomrule
    \end{tabular}
\end{table}

\begin{figure}[htbp]
    \centering
    \fbox{\rule{0pt}{52mm}\rule{0.84\linewidth}{0pt}}
    \caption{图位占位：CFG 引导实验结果图（待补）}
    \label{fig:algo-cfg-placeholder}
\end{figure}

\section{TTS 推理时搜索实验}

\subsection{实验设计}

TTS 的作用是把额外预算用于“候选生成与选择”，而非用于修改训练过程\cite{ttsdiffusion2025,ttssearch2025}。在本文中，TTS 被设计为推理层的外层搜索组件：先用当前采样配置生成多条候选轨迹，再根据目标一致性、前进性和平滑性综合打分，最终选择一条更适合执行的轨迹。实验重点比较不同预算下的轨迹质量收益，以及它与 DDIM、CFG 的组合关系。

\subsection{实验组织方式}

为了保证结论具有统计意义，本文将 TTS 实验组织为三层：单因素实验，用于比较是否开启 TTS 以及不同 budget；两两组合实验，用于比较 DDIM+TTS、CFG+TTS 的相互作用；三因素联合实验，用于比较 DDIM、CFG、TTS 同时开启时是否存在部署上更优的折中组合。

\subsection{结果占位与预期结论}

TTS 实验的核心问题不是“是否绝对提升所有指标”，而是“是否值得把 DDIM 节省下来的预算重新分配给候选搜索”。若实验表明少步 DDIM 结合中等预算 TTS 能在总延迟可接受的前提下显著改善目标一致性与轨迹稳定性，则说明三者具有互补性；若实验表明 CFG 与 TTS 同时增强会引入过强偏置，则说明二者需要更精细的配平。

\begin{table}[htbp]
    \centering
    \caption{TTS 推理时搜索实验结果占位}
    \label{tab:tts-results}
    \begin{tabular}{lcccc}
        \toprule
        配置 & Latency/ms & Forward Progress & Smoothness & 综合评分 \\
        \midrule
        No TTS & 待补 & 待补 & 待补 & 待补 \\
        TTS-8 & 待补 & 待补 & 待补 & 待补 \\
        TTS-16 & 待补 & 待补 & 待补 & 待补 \\
        TTS-32 & 待补 & 待补 & 待补 & 待补 \\
        \bottomrule
    \end{tabular}
\end{table}

\begin{figure}[htbp]
    \centering
    \fbox{\rule{0pt}{52mm}\rule{0.84\linewidth}{0pt}}
    \caption{图位占位：TTS 联合实验结果图（待补）}
    \label{fig:algo-tts-placeholder}
\end{figure}

\section{视觉编码器升级实验}

\subsection{实验目标与公平性问题}

视觉编码器升级实验的目标有两个：一是比较不同视觉编码器在统一框架下的潜力，二是为最终部署版训练筛选候选编码器。基于这一目的，本文将视觉编码器实验分为两个阶段：第一阶段使用统一 backbone-suite 框架进行公平筛选；第二阶段再对筛选出的候选做面向部署的最终训练。这样的两阶段设计能够同时满足“对比公平”和“训练结果可用于最终部署”两个需求。

\subsection{实验配置}

当前第一阶段已包含 EfficientNet-B0 baseline、backbone-suite 版 EfficientNet-B0、公平替换框架下的 DINOv2-Small、ConvNeXt-Tiny 和 ResNet-50。统一配置中，主要训练批大小为 32，训练轮数为 100。这样做的目标不是立即得出“谁绝对最好”，而是先判断哪些视觉编码器值得继续投入。

\subsection{已有结果与分析}

已有离线结果表明，DINOv2-Small 在部分单步代理指标上优于原始 EfficientNet-B0，说明更强视觉表征对局部条件建模具有潜力。但在闭环层面，原始 baseline 仍然表现最稳，ConvNeXt-Tiny 具有继续优化空间，而 DINOv2-Small 与 ResNet-50 目前尚不能写成“已经成功升级”的最终部署结论。其原因更可能来自当前统一替换框架与原始 NoMaD 结构之间的归纳偏置差异、冻结策略偏保守、以及距离预测头与编码器耦合导致的适配不足，而不能简单归因于“大模型编码器更差”。

因此，本文对视觉编码器实验的表述采用更谨慎的口径：第一阶段工作已经建立了统一比较框架，并识别出 DINOv2-Small 等候选方案；但“最终部署最优编码器”的结论仍需第二阶段面向部署的训练与闭环验证完成后才能正式写出。

\begin{table}[htbp]
    \centering
    \caption{视觉编码器升级实验结果占位}
    \label{tab:encoder-results}
    \begin{tabular}{lcccc}
        \toprule
        编码器 & Dist Pred Mean & Forward Progress & Smoothness & 闭环结论 \\
        \midrule
        EfficientNet-B0 & 待补 & 待补 & 待补 & 当前最稳定基线 \\
        EfficientNet-B0-Suite & 待补 & 待补 & 待补 & 公平对照待补 \\
        DINOv2-Small & 待补 & 待补 & 待补 & 第一阶段候选 \\
        ConvNeXt-Tiny & 待补 & 待补 & 待补 & 第一阶段候选 \\
        ResNet-50 & 待补 & 待补 & 待补 & 需继续检查 \\
        \bottomrule
    \end{tabular}
\end{table}

\begin{figure}[htbp]
    \centering
    \fbox{\rule{0pt}{52mm}\rule{0.84\linewidth}{0pt}}
    \caption{图位占位：视觉编码器升级实验结果图（待补）}
    \label{fig:algo-encoder-placeholder}
\end{figure}

\section{联合结论与进入闭环的候选方案}

综合本章离线实验，可得到三点阶段性结论。第一，在不改动训练结果的前提下，\texttt{ddim\_2} 是当前最值得进入闭环验证的加速配置。第二，CFG 的有效范围有限，更适合作为可调参数而不是默认固定设置。第三，视觉编码器升级工作已经完成第一阶段筛选框架搭建，但最终部署结论必须通过第二阶段训练和闭环验证给出。

从推理层组合关系看，DDIM、CFG 与 TTS 可以统一收口在同一推理模块内，但不能简单看作并列开关。更合理的组合方式是：先由 DDIM 决定采样调度，再由 CFG 决定条件引导强度，最后由 TTS 在候选轨迹层进行搜索与筛选。因此，本文建议后续闭环优先验证的主方案为“DDIM-2 + 可选 CFG + 可配置 TTS”，并以此作为第四章系统与仿真验证的输入。

\section{本章小结}

本章围绕 NoMaD 的算法层问题展开研究。首先复现并梳理了模型的输入输出形式、训练目标和统一推理接口；随后以基线实验为参照，对 DDIM 推理加速、CFG 引导增强、TTS 推理时搜索和视觉编码器升级进行了系统分析。当前结论表明，NoMaD 的部署潜力主要取决于推理层优化与视觉条件建模的共同效果，而不是单个指标的局部提升。下一章将在此基础上把这些候选方案接入统一闭环系统，并通过 MuJoCo 仿真验证它们在 Lite3 平台上的系统级可行性。
